% Gerado por scripts/migrate_markdown_to_latex.py.
% Fonte migrada: fases/01_validacao_conceitual/docs/conceitos.txt


Validação inicial dos conceitos e scripts do projeto



\section{Nota metodológica}






Os artigos nem sempre dão uma definição formal, no estilo 'X significa Y'. Muitas vezes eles usam o conceito dentro da arquitetura, da fórmula ou do procedimento experimental. Por isso, neste documento eu não tratei as definições como definições de dicionário. Tratei como definições operacionais:





o que este termo significa no nosso estudo -> por que essa interpretação vem da literatura -> qual comportamento precisa aparecer no código -> como isso foi testado.



A estrutura usada em cada item foi:




referência científica -> definição/conceito -> propriedade esperada no código -> script Python -> teste determinístico -> limite da conclusão




\section{1. Rede neural}


\begin{itemize}
\item Referência científica:
Goodfellow, Bengio e Courville, Deep Learning.
\end{itemize}

\begin{itemize}
\item Definição/conceito:
Neste estudo, 'rede neural' significa uma função computacional parametrizada,
composta por operações diferenciáveis organizadas em camadas. A justificativa é
que o livro trata redes neurais como modelos compostos por camadas e
parâmetros ajustáveis. No nosso caso, isso aparece como camadas lineares,
normalização e feed-forward.
\end{itemize}

\begin{itemize}
\item Propriedade esperada no código:
O bloco precisa receber um tensor, aplicar camadas com parâmetros e produzir
outro tensor como saída.
\end{itemize}

\begin{itemize}
\item Script Python:
Foi usada uma lógica de módulo neural com camadas lineares, normalização e
feed-forward.
\end{itemize}

\begin{itemize}
\item Teste determinístico:
Verificou-se se o bloco executa uma entrada fixa e gera saída com formato
esperado. Também foi usada uma rede Linear + ReLU + Linear como controle
negativo para mostrar que 'ser NN' não basta para ser Transformer.
\end{itemize}

\begin{itemize}
\item Limite da conclusão:
Isso valida que existe uma NN no experimento, mas não valida que ela seja uma
Transformer.
\end{itemize}



\section{2. Deep learning}


\begin{itemize}
\item Referência científica:
Goodfellow, Bengio e Courville, Deep Learning.
\end{itemize}

\begin{itemize}
\item Definição/conceito:
Neste estudo, 'deep learning' não é um resultado que estamos tentando provar.
Ele serve como contexto: usamos uma rede com várias transformações encadeadas,
e não uma única fórmula isolada. A justificativa vem do uso de múltiplas
camadas para compor representações.
\end{itemize}

\begin{itemize}
\item Propriedade esperada no código:
Deve haver uma sequência de transformações: projeções, attention, projeção de
saída, normalização e feed-forward.
\end{itemize}

\begin{itemize}
\item Script Python:
Foi usada uma lógica de bloco neural composto por mais de uma etapa
computacional.
\end{itemize}

\begin{itemize}
\item Teste determinístico:
Verificou-se se o bloco completo recebe a entrada, atravessa essas etapas e
mantém o formato esperado da sequência.
\end{itemize}

\begin{itemize}
\item Limite da conclusão:
Isso não valida aprendizado, treinamento ou generalização. Só valida a
estrutura computacional usada no recorte.
\end{itemize}



\section{3. Transformer}


\begin{itemize}
\item Referência científica:
Vaswani et al., Attention Is All You Need, 2017.
\end{itemize}

\begin{itemize}
\item Definição/conceito:
Neste estudo, 'Transformer' é entendido operacionalmente como uma arquitetura
para sequências baseada em attention, sem recorrência ou convolução como
mecanismo central. A justificativa vem do artigo de Vaswani, que apresenta o
Transformer como arquitetura baseada em mecanismos de atenção e descreve blocos
com self-attention, feed-forward, residual, normalização e informação
posicional.
\end{itemize}

\begin{itemize}
\item Propriedade esperada no código:
Para pertencer à família Transformer, o bloco precisa ter entrada sequencial,
Q/K/V, scaled dot-product attention, dependência entre posições, informação
posicional e componentes de bloco.
\end{itemize}

\begin{itemize}
\item Script Python:
Foi usada uma lógica de auditoria que procura essas propriedades no bloco e
classifica o resultado.
\end{itemize}

\begin{itemize}
\item Teste determinístico:
O bloco foi testado contra a checklist arquitetural e classificado como bloco
Transformer simplificado. Uma NN comum foi rejeitada.
\end{itemize}

\begin{itemize}
\item Limite da conclusão:
Isso sustenta dizer que usamos um bloco da família Transformer. Não sustenta
dizer que implementamos a Transformer completa do artigo.
\end{itemize}



\section{4. Bloco Transformer simplificado}


\begin{itemize}
\item Referência científica:
Vaswani et al., Attention Is All You Need, 2017.
\end{itemize}

\begin{itemize}
\item Definição/conceito:
Neste estudo, 'bloco Transformer simplificado' significa um recorte controlado
com as partes essenciais de um bloco: self-attention multi-head, projeção de
saída, residual, normalização, feed-forward e informação posicional. A
justificativa é que o artigo descreve cada camada do encoder como composta por
multi-head self-attention e uma rede feed-forward, com residual e layer norm.
\end{itemize}

\begin{itemize}
\item Propriedade esperada no código:
O bloco precisa combinar sequência de entrada, Q/K/V, attention, recombinação
de heads, projeção de saída, residual, normalização e FFN.
\end{itemize}

\begin{itemize}
\item Script Python:
Foi usada uma lógica de bloco único, pequeno e controlado, em vez de uma pilha
completa de encoder/decoder.
\end{itemize}

\begin{itemize}
\item Teste determinístico:
Verificou-se shape de entrada/saída, presença dos componentes e aprovação na
auditoria arquitetural.
\end{itemize}

\begin{itemize}
\item Limite da conclusão:
O bloco é adequado para estudar trechos centrais da Transformer. Ele não é a
arquitetura completa.
\end{itemize}



\section{5. Projeção linear densa}


\begin{itemize}
\item Referência científica:
Goodfellow, Bengio e Courville, Deep Learning; Vaswani et al., 2017.
\end{itemize}

\begin{itemize}
\item Definição/conceito:
Neste estudo, projeção linear densa é a operação Y = XW\textasciicircum{}T + b. A justificativa
é que Vaswani usa projeções lineares aprendidas para gerar Q, K, V e também
para projetar a saída da multi-head attention.
\end{itemize}

\begin{itemize}
\item Propriedade esperada no código:
A saída deve ser exatamente o produto matricial da entrada pelos pesos
transpostos, com soma de viés quando houver.
\end{itemize}

\begin{itemize}
\item Script Python:
Foi usada uma lógica direta de multiplicação matricial e soma opcional de viés.
\end{itemize}

\begin{itemize}
\item Teste determinístico:
Comparou-se um caso pequeno com resultado manual conhecido e com cálculo NumPy
independente.
\end{itemize}

\begin{itemize}
\item Limite da conclusão:
Isso valida a operação usada para projeção. Não valida que os pesos sejam
treinados ou semanticamente úteis.
\end{itemize}



\section{6. Q, K e V}


\begin{itemize}
\item Referência científica:
Vaswani et al., Attention Is All You Need, 2017.
\end{itemize}

\begin{itemize}
\item Definição/conceito:
Neste estudo, Q, K e V significam três projeções lineares da representação de
entrada usadas pela attention. A justificativa é que Vaswani descreve queries,
keys e values como matrizes projetadas linearmente antes da attention.
\end{itemize}

\begin{itemize}
\item Propriedade esperada no código:
A mesma sequência precisa gerar três tensores separados: query, key e value.
\end{itemize}

\begin{itemize}
\item Script Python:
Foi usada uma lógica de três projeções lineares independentes aplicadas antes
da attention.
\end{itemize}

\begin{itemize}
\item Teste determinístico:
A auditoria verificou se existem as três projeções e se o núcleo de attention
recebe Q, K e V.
\end{itemize}

\begin{itemize}
\item Limite da conclusão:
Isso valida a estrutura Q/K/V do bloco. Não valida encoder-decoder completo.
\end{itemize}



\section{7. Scaled dot-product attention}


\begin{itemize}
\item Referência científica:
Vaswani et al., Attention Is All You Need, 2017; documentação PyTorch da
scaled\_\allowbreak{}dot\_\allowbreak{}product\_\allowbreak{}attention.
\end{itemize}

\begin{itemize}
\item Definição/conceito:
Neste estudo, scaled dot-product attention significa exatamente a operação
softmax(QK\textasciicircum{}T / sqrt(d\_\allowbreak{}k))V. A justificativa é que essa é a fórmula apresentada
por Vaswani para a attention escalonada e é a mesma lógica documentada na API
de referência do PyTorch.
\end{itemize}

\begin{itemize}
\item Propriedade esperada no código:
O código precisa calcular QK\textasciicircum{}T, dividir por sqrt(d\_\allowbreak{}k), aplicar softmax e usar
esses pesos para combinar V.
\end{itemize}

\begin{itemize}
\item Script Python:
Foi usada a lógica literal da fórmula, sem otimização escondida.
\end{itemize}

\begin{itemize}
\item Teste determinístico:
Comparou-se a saída com um valor esperado pequeno, com implementação NumPy
independente e com a função de referência do PyTorch.
\end{itemize}

\begin{itemize}
\item Limite da conclusão:
Isso valida o núcleo matemático da attention. Attention isolada não é
Transformer completa.
\end{itemize}



\section{8. Self-attention}


\begin{itemize}
\item Referência científica:
Vaswani et al., Attention Is All You Need, 2017.
\end{itemize}

\begin{itemize}
\item Definição/conceito:
Neste estudo, self-attention significa attention em que Q, K e V vêm da mesma
sequência. A justificativa é que Vaswani descreve self-attention no encoder
como o caso em que queries, keys e values vêm da saída da camada anterior do
próprio encoder.
\end{itemize}

\begin{itemize}
\item Propriedade esperada no código:
A alteração de uma posição da sequência deve poder afetar a saída de outra
posição.
\end{itemize}

\begin{itemize}
\item Script Python:
Foi usada uma lógica em que a própria sequência alimenta Q, K e V.
\end{itemize}

\begin{itemize}
\item Teste determinístico:
Alterou-se um token e verificou-se se outra posição da saída mudava.
\end{itemize}

\begin{itemize}
\item Limite da conclusão:
Isso valida dependência entre posições. Não valida qualidade linguística ou
desempenho em tarefa real.
\end{itemize}



\section{9. Multi-head attention}


\begin{itemize}
\item Referência científica:
Vaswani et al., Attention Is All You Need, 2017.
\end{itemize}

\begin{itemize}
\item Definição/conceito:
Neste estudo, multi-head attention significa dividir a representação em
múltiplas cabeças, aplicar scaled dot-product attention em paralelo e
recombinar os resultados. A justificativa é que Vaswani descreve a projeção de
Q/K/V em várias heads, seguida de attention paralela, concatenação e nova
projeção.
\end{itemize}

\begin{itemize}
\item Propriedade esperada no código:
O tensor precisa ser separado em heads, processado por attention e reconstruído
no formato original.
\end{itemize}

\begin{itemize}
\item Script Python:
Foi usada uma lógica de split de heads, attention por head e recombinação.
\end{itemize}

\begin{itemize}
\item Teste determinístico:
Comparou-se o resultado com a attention do PyTorch aplicada em cada head e
recombinada manualmente.
\end{itemize}

\begin{itemize}
\item Limite da conclusão:
Isso valida a mecânica multi-head. Não prova que a quantidade de heads
escolhida seja ótima.
\end{itemize}



\section{10. Informação posicional}


\begin{itemize}
\item Referência científica:
Vaswani et al., Attention Is All You Need, 2017.
\end{itemize}

\begin{itemize}
\item Definição/conceito:
Neste estudo, informação posicional significa um sinal somado à sequência para
representar ordem. A justificativa é que Vaswani explica que, como o modelo não
usa recorrência nem convolução, precisa injetar informação sobre posição dos
tokens.
\end{itemize}

\begin{itemize}
\item Propriedade esperada no código:
A entrada sequencial precisa receber informação dependente da posição antes da
attention.
\end{itemize}

\begin{itemize}
\item Script Python:
Foi usada uma lógica de positional encoding sinusoidal somado à sequência.
\end{itemize}

\begin{itemize}
\item Teste determinístico:
A auditoria verificou a existência de informação posicional antes das
projeções.
\end{itemize}

\begin{itemize}
\item Limite da conclusão:
Isso valida a presença de posição. Não prova que esse encoding seja o melhor
possível.
\end{itemize}



\section{11. Inferência}


\begin{itemize}
\item Referência científica:
Documentação PyTorch sobre modo de avaliação e reprodutibilidade.
\end{itemize}

\begin{itemize}
\item Definição/conceito:
Neste estudo, inferência significa executar o bloco para obter saída sem
treinar, sem atualizar pesos e com comportamento determinístico. Essa definição
é operacional porque os benchmarks medirão execução, não aprendizado.
\end{itemize}

\begin{itemize}
\item Propriedade esperada no código:
A mesma entrada, no mesmo estado do modelo, deve produzir a mesma saída.
\end{itemize}

\begin{itemize}
\item Script Python:
Foi usada uma lógica de execução em modo de avaliação e sem atualização de
parâmetros.
\end{itemize}

\begin{itemize}
\item Teste determinístico:
A mesma entrada foi executada duas vezes e as saídas foram comparadas.
\end{itemize}

\begin{itemize}
\item Limite da conclusão:
Isso valida determinismo do núcleo. Não substitui repetição estatística nos
benchmarks.
\end{itemize}



\section{12. KV cache}


\begin{itemize}
\item Referência científica:
Literatura de inferência autoregressiva em Transformers e otimizações de KV
cache.
\end{itemize}

\begin{itemize}
\item Definição/conceito:
Neste estudo, KV cache significa armazenar chaves e valores já calculados para
não recomputá-los a cada novo token. A definição é operacional porque o que
importa para nós é reaproveitamento de K/V e equivalência de saída.
\end{itemize}

\begin{itemize}
\item Propriedade esperada no código:
A saída com cache deve bater com a saída sem cache, usando menos recomputação
de K/V.
\end{itemize}

\begin{itemize}
\item Script Python:
Foi usada uma lógica token a token que acumula K/V anteriores e reutiliza esse
cache na attention causal.
\end{itemize}

\begin{itemize}
\item Teste determinístico:
Comparou-se a saída com cache contra a saída causal completa e contou-se o
reaproveitamento de K/V.
\end{itemize}

\begin{itemize}
\item Limite da conclusão:
Isso valida equivalência algorítmica. Não prova ganho de hardware sem medição.
\end{itemize}



\section{13. Pruning}


\begin{itemize}
\item Referência científica:
Tang et al., A Survey on Transformer Compression, 2024; documentação PyTorch
pruning.
\end{itemize}

\begin{itemize}
\item Definição/conceito:
Neste estudo, pruning significa remover ou zerar pesos segundo algum critério,
como magnitude. A justificativa é que a literatura de compressão usa pruning
como técnica de reduzir parâmetros/conexões, mas o ganho físico só existe se a
execução conseguir explorar essa remoção.
\end{itemize}

\begin{itemize}
\item Propriedade esperada no código:
Após pruning, deve existir máscara ou aumento verificável de pesos zerados.
\end{itemize}

\begin{itemize}
\item Script Python:
Foi usada uma lógica de pruning por magnitude em camada linear pequena.
\end{itemize}

\begin{itemize}
\item Teste determinístico:
Mediu-se a quantidade de zeros e verificou-se a máscara gerada.
\end{itemize}

\begin{itemize}
\item Limite da conclusão:
Isso valida pruning conceitual/algorítmico. Não valida ganho de hardware se a
matriz continuar sendo executada como densa.
\end{itemize}



\section{14. Sparsity}


\begin{itemize}
\item Referência científica:
Literatura de pruning e compressão de redes neurais.
\end{itemize}

\begin{itemize}
\item Definição/conceito:
Neste estudo, sparsity significa a razão entre elementos iguais a zero e o
total de elementos. A justificativa é que essa é a propriedade observável que
indica se o pruning gerou esparsidade.
\end{itemize}

\begin{itemize}
\item Propriedade esperada no código:
O código deve contar zeros e dividir pelo número total de elementos.
\end{itemize}

\begin{itemize}
\item Script Python:
Foi usada uma lógica direta de contagem de zeros.
\end{itemize}

\begin{itemize}
\item Teste determinístico:
Comparou-se a taxa calculada com tensor pequeno de esparsidade conhecida.
\end{itemize}

\begin{itemize}
\item Limite da conclusão:
Sparsity mede zeros. Não mede automaticamente redução de tempo ou memória real.
\end{itemize}



\section{15. Quantização}


\begin{itemize}
\item Referência científica:
Tang et al., A Survey on Transformer Compression, 2024; TurboQuant;
documentação torchao/PyTorch.
\end{itemize}

\begin{itemize}
\item Definição/conceito:
Neste estudo, quantização significa representar valores com menor precisão,
especialmente int8, mantendo uma escala para aproximar os valores originais. A
justificativa é que a literatura de compressão trata quantização como mudança
de representação numérica para reduzir custo/armazenamento, com erro medido em
relação ao baseline.
\end{itemize}

\begin{itemize}
\item Propriedade esperada no código:
Deve aparecer dtype menor, escala de quantização, armazenamento menor ou
estimado e erro numérico em relação ao tensor original.
\end{itemize}

\begin{itemize}
\item Script Python:
Foi usada uma lógica de quantização simétrica int8, seguida de dequantização
apenas para medir erro.
\end{itemize}

\begin{itemize}
\item Teste determinístico:
Verificou-se dtype int8, bytes por elemento, erro após dequantização e uma
configuração int8 weight-only com torchao.
\end{itemize}

\begin{itemize}
\item Limite da conclusão:
Isso valida quantização numérica/algorítmica. Ganho de hardware exige kernel
ou caminho de execução compatível.
\end{itemize}



\section{16. Latência}


\begin{itemize}
\item Referência científica:
Literatura de desempenho computacional e documentação PyTorch/CUDA.
\end{itemize}

\begin{itemize}
\item Definição/conceito:
Neste estudo, latência significa tempo por execução do bloco em um cenário
definido. A definição é operacional porque será medida diretamente no
benchmark.
\end{itemize}

\begin{itemize}
\item Propriedade esperada no código:
A medição precisa usar warmup, repetições e sincronização quando o dispositivo
for CUDA.
\end{itemize}

\begin{itemize}
\item Script Python:
Foi usada uma lógica de aquecimento, medição repetida e sincronização de GPU.
\end{itemize}

\begin{itemize}
\item Teste determinístico:
Verificou-se se o protocolo registra média, p50 e p95 de latência no schema de
saída.
\end{itemize}

\begin{itemize}
\item Limite da conclusão:
O protocolo está pronto. A conclusão de desempenho só vem após benchmark.
\end{itemize}



\section{17. Throughput}


\begin{itemize}
\item Referência científica:
Literatura de desempenho computacional.
\end{itemize}

\begin{itemize}
\item Definição/conceito:
Neste estudo, throughput significa tokens processados por segundo. A definição
é operacional porque depende diretamente de batch, sequência e tempo medido.
\end{itemize}

\begin{itemize}
\item Propriedade esperada no código:
Deve calcular tokens/segundo a partir do número de tokens processados e da
latência.
\end{itemize}

\begin{itemize}
\item Script Python:
Foi usada uma lógica de batch\_\allowbreak{}size * seq\_\allowbreak{}len dividido pelo tempo médio.
\end{itemize}

\begin{itemize}
\item Teste determinístico:
Verificou-se se o schema dos resultados exige registro de throughput.
\end{itemize}

\begin{itemize}
\item Limite da conclusão:
O cálculo está definido. A comparação depende da execução experimental.
\end{itemize}



\section{18. Memória}


\begin{itemize}
\item Referência científica:
Literatura de arquitetura de computadores e documentação PyTorch.
\end{itemize}

\begin{itemize}
\item Definição/conceito:
Neste estudo, memória significa armazenamento/alocação associado ao cenário
executado. A definição é operacional porque será registrada como pico medido ou
estimativa teórica.
\end{itemize}

\begin{itemize}
\item Propriedade esperada no código:
Deve haver uma medida ou estimativa comparável entre baseline, pruning e
quantização.
\end{itemize}

\begin{itemize}
\item Script Python:
Foi usada uma lógica de registrar pico de memória quando disponível e estimar
armazenamento dos tensores.
\end{itemize}

\begin{itemize}
\item Teste determinístico:
Verificou-se se o schema exige coluna de memória.
\end{itemize}

\begin{itemize}
\item Limite da conclusão:
Menor armazenamento teórico não prova automaticamente menor uso real de
hardware.
\end{itemize}



\section{Conclusão geral}






O documento agora separa definição operacional de definição formal. A ideia não é fingir que todo artigo traz uma frase fechada definindo cada termo. A ideia é mostrar que, quando a literatura usa um conceito, nós extraímos dele uma propriedade testável e verificamos essa propriedade no código.



Com isso, considero validado neste momento:

\begin{itemize}
\item que o núcleo experimental implementa trechos centrais do Transformer de
Vaswani: Q/K/V, scaled dot-product attention, self-attention, multi-head,
informação posicional, residual, normalização e FFN;
\item que o núcleo deve ser chamado de bloco Transformer simplificado, não de
Transformer completa;
\item que pruning, sparsity, quantização, latência, throughput e memória foram
definidos de forma operacional e testável;
\item que ganho real de hardware ainda não foi afirmado, porque isso depende dos
benchmarks controlados.
\end{itemize}
